% Contenido científico íntegro normalizado del frontmatter y cortes
% Residencia material del ensamblaje sucesor de 102 capítulos.
% parte_i.tex:1--419,529--659 + parte_iii.tex:1--908.
% Se excluyen sólo la portadilla autónoma y sus órdenes de página.
\part{Correspondances radiales généalogiques sur l'état enrichi : domaine, codomaine et conservation de l'histoire hélicoïdale}
Cette publication développe une théorie généalogique des correspondances spirales à partir de la composition \(APP\to\TRIT\to\TPK\). L'objet primaire n'est pas une courbe présupposée dans \(\R^2\), mais une histoire compatible de l'état enrichi, avec ses feuillets arithmétiques, son orientation, sa retenue, son bord, sa voie et sa mémoire. La courbe apparaît à la fin comme publication archimédienne. En conséquence, deux figures égales dans \(\R^2\) peuvent provenir d'histoires distinctes, et un cercle visible peut posséder une mémoire radiale intrinsèquement nulle ou être l'ombre d'une hélice dont la mémoire a été éliminée par la projection.

L'ouvrage conserve et développe les résultats exacts du corpus HMT sur la courbure mixte d'APP, les trois régimes quadratiques de TRIT, l'holonomie nonadique \(\Gamma_9\), le retour de phase avec avancée de mémoire, l'hélice nonadique, l'extension dodécaphasique non scindée et la vis spectrale de Pell. Sur ces résultats, il construit une formalisation nouvelle : l'enveloppe affine universelle, le lecteur radial enrichi par préfixes, sa naturalité sous troncature, son équivariance avec \(\Gamma_9\) et sa publication puissance--logarithmique.

La construction n'altère pas la théorie des masses et ne rouvre pas l'électron central. L'unique point fixe \((5,5)\) de l'atlas matériel est utilisé uniquement comme ancrage ultérieur d'une interface ; jamais comme sélecteur rétrospectif de la géométrie radiale.

\medskip
\begin{statusbox}[title={Stratification probatoire par domaine et codomaine}]
L'exposé distingue les théorèmes internes précédemment établis, les théorèmes démontrés dans cet ouvrage, les certifications finies et les réalisations physiques. Un résultat algébrique interne est pleinement affirmé dans son domaine. Toute réalisation physique ultérieure conserve ses prémisses et n'intervient pas dans la sélection du générateur HMT.
\end{statusbox}

\part{Construction et résultats principaux du lecteur radial enrichi}
On construit une théorie des correspondances spirales dont l'antécédent est l'état enrichi de l'holographie modulaire triadique et dont la sortie est une famille de courbes dans \(\R^2\) ou \(\R^3\). APP fournit un support bidimensionnel avec deux évaluations non simultanément compatibles, \(\Sigma\) et \(\Pi\) ; leur composition mixte possède une courbure discrète orientée \(\pm1\). TRIT type les régimes elliptique, parabolique et hyperbolique sans les confondre avec le caractère radial ni avec la courbure de Frenet. TPK sélectionne, transporte, mémorise, actualise et ramène l'état complet. L'holonomie \(\Gamma_9\) ramène la phase, mais accroît la profondeur de mémoire et transporte le bord.

Le lecteur initialement suggéré vers \((p,a,\Lambda,C,m,r,\theta)\) est corrigé de manière décisive. Les coordonnées \((r,\theta)\) appartiennent à la publication archimédienne, et la paire finale \((\Lambda,C)\) ne conserve pas ses préfixes. On définit donc un lecteur interne
\[
 \mathcal P^{\enr}_{\rad,N}:\Hist_N^{\enr}(\TPK)
 \longrightarrow \mathcal D^{\enr}_{\rad,N}
\]
qui enregistre, à chaque profondeur, les cocycles de mode et d'orientation, le facteur affine local, l'holonomie préfixe et le registre complet. On démontre sa naturalité sous troncature, son extension à la limite inverse, son équivariance avec \(\Gamma_9\), sa conservation de la retenue, du bord et de la mémoire, ainsi que sa compatibilité avec l'inversion des voies dans le secteur réversible.

Après la construction du continu, on applique un caractère radial et l'on introduit la coordonnée puissance--logarithmique
\[
 L_p(x)=\frac{x^p-1}{p},\qquad L_0(x)=\log x.
\]
Sa loi de composition rend visible l'articulation somme--produit d'APP. La translation uniforme dans \(L_p\) publie, pour \(p=2,1,0,-1,-2\), les spirales de Fermat, d'Archimède, logarithmique, hyperbolique et lituus. L'holonomie affine générale produit une classification plus large au moyen de \((p,[\Lambda,C])\). La vis spectrale \(V^{-1}\mathscr SV=(2+\sqrt3)\ii\mathscr S\) apparaît comme cas exact : sa projection transverse est une spirale logarithmique et son relèvement conserve la mémoire hélicoïdale.

La publication s'achève sur la géométrie différentielle des courbes publiées, une réalisation électromagnétique restreinte de l'hélice phase--position, l'interface avec l'électron central et un système de réfutateurs qui empêche de confondre apparence, généalogie et réalisation physique.

\bigskip
\noindent\textbf{Mots-clés :} APP ; TRIT ; Topological Phase Kernel ; holonomie nonadique ; cocycle affine ; spirales ; suspension hélicoïdale ; mémoire ; double projection ; continu ; mécanique dimensionnelle.

\part{Résultats directeurs de la publication radiale APP--TRIT--TPK}
\begin{enumerate}[label=\textbf{R\arabic*.},leftmargin=12mm]
\item La géométrie visible est une publication non fidèle de l'état enrichi : l'égalité dans \(\R^2\) n'implique pas l'égalité de généalogie.
\item La courbure mixte d'APP distingue l'ordre somme--produit avant toute courbe réelle.
\item \(\tau\), \(p\) et la courbure de Frenet sont des invariants de niveaux distincts et ne peuvent être identifiés.
\item L'holonomie nonadique ramène la phase, fait avancer la mémoire et transporte le bord ; la projection angulaire périodique admet donc un relèvement hélicoïdal nonadique qui conserve la profondeur du retour.
\item L'information radiale complète exige de conserver tous les facteurs et toutes les holonomies préfixes ; la paire finale \((\Lambda,C)\) est insuffisante.
\item Le lecteur radial enrichi est naturel sous troncature et équivariant sous \(\Gamma_9\).
\item L'involution bidirectionnelle produit des voies conjuguées d'ascension et de descente dans le secteur réversible.
\item La famille \(L_p\) publie les cinq classes puissance--logarithmiques comme sorties d'une même loi, sans les choisir à l'avance.
\item La vis spectrale de Pell constitue une réalisation exacte de rotation et d'auto-échelle dont la projection est logarithmique.
\item Un cercle peut être un état radialement immobile ou l'ombre d'une hélice à registre non nul.
\item L'électron \((5,5)\) demeure le centre matériel démontré ; la signature spirale est ajoutée comme lecteur, non comme substitut de l'atlas \(81/56/13/324\).
\item Les réalisations électromagnétiques et matérielles sont postérieures à la construction HMT et conservent leur propre typage physique.
\end{enumerate}

\part{Notation typée et niveaux causaux de la publication radiale}
\begin{longtable}{>{\raggedright\arraybackslash}p{.23\textwidth}p{.67\textwidth}}
\toprule
Symbole & Signification\\
\midrule
\endfirsthead
\toprule
Symbole & Signification\\
\midrule
\endhead
\(X_9\) & Tore nonadique \((\Z/9\Z)^2\), support bidimensionnel d'APP.\\
\(\Sigma,\Pi\) & Évaluations additive et multiplicative du même support.\\
\(\tau\in\{-1,0,+1\}\) & Type local TRIT de transport : hyperbolique, parabolique ou elliptique selon la convention quadratique déclarée.\\
\(U_t\) & Transition typée \(\mathrm{Upd}_t\circ\mathrm{Tra}_t\circ\mathrm{Sel}_t\).\\
\(\Gamma_9\) & Holonomie de neuf transitions sur l'état enrichi.\\
\(p\) & Cocycle ou degré radial après définition de son extraction interne ; ce n'est pas \(\tau\).\\
\(a\) & Caractère orienté local \(a=\tau_+-\tau_-\) ; il se distingue du degré radial \(p=\tau_++\tau_-\).\\
\((\Lambda,C)\) & Endomorphisme affine radial : partie linéaire et translation.\\
\(H_j\) & Holonomie affine préfixe jusqu'à la profondeur \(j\).\\
\(L_p,\exp_p\) & Coordonnée puissance--logarithmique et son inverse.\\
\(\kappa_{\mathrm{F}},\mathcal T_{\mathrm{F}}\) & Courbure et torsion de Frenet de la courbe publiée.\\
\(\mathcal P^{\enr}_{\rad}\) & Lecteur généalogique radial interne.\\
\(\ev_{\arq}^{\mathrm{esp}}\) & Publication archimédienne de la signature spirale.\\
\bottomrule
\end{longtable}

\begin{resultbox}[title={Composition causale d'APP--TRIT--TPK à l'évaluation archimédienne radiale}]
\[
\APP\longrightarrow\TRIT\longrightarrow\TPK
\longrightarrow X_{\TPK}^{\enr}
\longrightarrow\mathcal C_{\mathrm{cont}}^{\mathrm{disc}}
\longrightarrow\mathcal P_{
  \rad}^{\enr}
\longrightarrow\ev_{\arq}^{\mathrm{esp}}
\longrightarrow\bigl(\R^2\ \text{ou}\ \R^3\bigr).
\]
La courbe est une sortie de la chaîne générative ; elle n'intervient dans la sélection ni de l'état ni de ses coefficients.
\end{resultbox}
\part{Fondement généalogique de la géométrie radiale}

\chapter{La géométrie visible comme publication de l'état enrichi}

\section{L'inversion de l'ordre explicatif}

La géométrie conventionnelle commence par un ensemble continu et y définit des courbes, des paramètres et des champs. La construction développée ici inverse cet ordre. On produit d'abord une histoire compatible sur un support fini, puis on conserve ses données de transport et ce n'est qu'à la fin que l'on publie une coordonnée réelle. Le schéma causal est
\begin{equation}
 \APP\longrightarrow\TRIT\longrightarrow\TPK
 \longrightarrow X_{\TPK}^{\enr}
 \longrightarrow\mathcal C_{\mathrm{cont}}^{\mathrm{disc}}
 \longrightarrow\ev_{\arq}
 \longrightarrow\R^d.
 \label{espiral:eq:cadena-genealogica-espiral}
\end{equation}
Il ne s'agit pas d'orner une courbe classique d'étiquettes discrètes. L'histoire est l'objet ; la courbe est l'une de ses lectures.

\begin{definition}[Génération géométrique avec traçabilité]
Une construction géométrique est \emph{constructive--générative} lorsqu'elle :
\begin{enumerate}
\item produit son objet à partir d'opérations déclarées sur un domaine antérieur à la géométrie réelle ;
\item conserve l'ordre des opérations et les coordonnées nécessaires pour le reconstruire ;
\item obtient l'apparence géométrique au moyen d'une application ultérieure ;
\item distingue l'égalité des apparences de l'équivalence des histoires.
\end{enumerate}
\end{definition}

Ce critère est plus fort qu'une paramétrisation. Une paramétrisation \(t\mapsto\gamma(t)\) présuppose la courbe et change son mode de parcours. Une généalogie \(x=(x_n)_{n\ge0}\), en revanche, spécifie les troncatures finies, les transitions admissibles, les feuillets arithmétiques, la retenue et la mémoire qui rendent possible la sortie \(\gamma\).

\section{État, histoire et publication}

Soit \(X_{\TPK}^{\enr}\) l'espace des états enrichis. Une histoire finie de longueur \(N\) est
\[
 \gamma_N=(x_0\xrightarrow{e_1}x_1\xrightarrow{e_2}\cdots
 \xrightarrow{e_N}x_N),
\]
où chaque flèche est compatible avec la transition TPK typée. Nous associons à chaque préfixe son registre
\[
 \Led_j(\gamma_N)=
 (R_j,Q_j,\kappa_j,\sigma_j,o_j,g_j,C_j,B_j,
  \lambda_j,m_j,s_j),
\]
qui enregistre, selon l'emplacement concret, le résidu, le quotient, la retenue, le feuillet, l'orientation, la phase, le cylindre, le bord, la voie, la mémoire et la survie. Le contenu exact peut être enrichi, mais ne peut être réduit si la coordonnée éliminée intervient dans une continuation.

La publication réelle est une composition
\[
 \Hist_N^{\enr}(\TPK)
 \xrightarrow{\ \mathcal P_N^{\enr}\ }
 \mathcal D_N^{\enr}
 \xrightarrow{\ \ev_N\ }
 Y,
\]
avec \(Y=\R,\R^2\) ou \(\R^3\). La première application conserve la généalogie ; la seconde publie. Si \(\ev_N(\gamma)=\ev_N(\eta)\), on conclut seulement que les deux histoires appartiennent à la même fibre d'évaluation.

\begin{proposition}[Non-fidélité de la géométrie visible]
Si l'application d'évaluation élimine au moins une coordonnée du registre qui admet deux valeurs compatibles, il existe alors des histoires distinctes ayant la même sortie géométrique.
\end{proposition}
\begin{proof}
Soit \(c\) une coordonnée éliminée et soient \(\gamma,\eta\) deux histoires compatibles qui coïncident en toutes les coordonnées conservées et diffèrent en \(c\). Par définition de l'évaluation, \(\ev(\gamma)=\ev(\eta)\), tandis que \(\gamma\neq\eta\) car leurs registres diffèrent. Par conséquent, \(\ev\) n'est pas injective.
\end{proof}

La proposition n'est pas une insuffisance de la géométrie réelle. Elle exprime sa fonction : contracter une structure plus riche en une coordonnée utilisable. La thèse centrale consiste précisément à ne pas confondre cette contraction avec l'objet contracté.

\section{Cercle intrinsèque et cercle projeté}

Considérons l'application
\[
 c_0(t)=(R\cos t,R\sin t),
\]
et l'hélice
\[
 h(t)=(R\cos t,R\sin t,bt).
\]
La projection \(\pi_{xy}(x,y,z)=(x,y)\) satisfait \(\pi_{xy}\circ h=c_0\). Le même cercle peut représenter :
\begin{enumerate}
\item un état dont l'incrément radial et la mémoire sont nuls ;
\item l'ombre d'un état d'incrément radial nul et de mémoire axiale \(b\neq0\).
\end{enumerate}
La forme visible ne tranche pas entre les deux cas. Le registre est nécessaire.

\begin{figure}[htbp]
\centering
\begin{subfigure}[t]{.46\textwidth}
\centering
\begin{tikzpicture}[scale=.82]
  \draw[->,HMTGray] (-2.2,0)--(2.2,0);
  \draw[->,HMTGray] (0,-2.2)--(0,2.2);
  \draw[HMTBlue,very thick] (0,0) circle (1.6);
  \fill[HMTNavy] (1.6,0) circle (1.7pt);
  \node[below,HMTGray,font=\small\sffamily] at (0,-2.25) {mémoire visible nulle};
\end{tikzpicture}
\caption{Cercle comme objet plan.}
\end{subfigure}\hfill
\begin{subfigure}[t]{.46\textwidth}
\centering
\begin{tikzpicture}[x={(.9cm,0cm)},y={(.32cm,.14cm)},z={(0cm,.72cm)},scale=.74]
  \draw[->,HMTGray] (0,0,0)--(0,0,6.8);
  \draw[HMTBlue,very thick,domain=0:720,samples=160,smooth,variable=\t]
    plot ({1.45*cos(\t)},{1.45*sin(\t)},{\t/110});
  \draw[HMTRed,dashed,thick] (0,0,0) ellipse (1.45 and .47);
  \node[right,HMTGray,font=\small\sffamily] at (0,0,6.8) {mémoire};
\end{tikzpicture}
\caption{La même ombre, avec mémoire relevée.}
\end{subfigure}
\caption{La géométrie visible ne reconstruit pas à elle seule la généalogie.}
\label{espiral:fig:circulo-helice}
\end{figure}

\begin{resultbox}[title={Défaut de fidélité de l'évaluation archimédienne sur les classes généalogiques}]
Un cercle peut posséder une mémoire intrinsèquement nulle et peut aussi être la publication d'une histoire hélicoïdale. L'égalité de la courbe n'autorise pas à identifier les états qui la produisent.
\end{resultbox}

\chapter{APP : support bidimensionnel, feuillets arithmétiques et courbure mixte}

\section{La cellule nonadique et la naissance bidimensionnelle}

La cellule visible
\[
 E_9=\{1,2,3,4,5,6,7,8,9\}
\]
porte simultanément une phase et une carte de représentants positifs. Pour \(n\ge1\),
\[
 \rho_9(n)=1+((n-1)\bmod9),\qquad
 q_9(n)=\left\lfloor\frac{n-1}{9}\right\rfloor,
\]
de sorte que
\[
 n=9q_9(n)+\rho_9(n).
\]
Le produit cartésien de deux cellules construit
\[
 X_9=(\Z/9\Z)^2.
\]
Ce n'est pas une droite répétée, mais un support véritablement bidimensionnel. La complétion \(9\times9\) peut être lue comme un noyau \(8\times8\) et un gnomon de \(17\) cellules :
\[
81=64+17.
\]

Sur chaque \((i,j)\in X_9\) agissent deux évaluations du même support :
\[
 \Sigma(i,j)=\operatorname{dr}(i+j),\qquad
 \Pi(i,j)=\operatorname{dr}(ij),
\]
où \(\operatorname{dr}\) publie le représentant nonadique et le relèvement conserve le quotient. APP n'additionne pas deux espaces ; il conserve un seul espace avec deux lectures.

\section{Résidu, quotient et retenue}

Les identités relevées
\[
 i+j=R^+_{ij}+9q^+_{ij},\qquad
 ij=R^\times_{ij}+9q^\times_{ij}
\]
distinguent la sortie visible \(R\) du relèvement \(q\). Si seul le résidu est conservé, des opérations distinctes peuvent se confondre dans la même cellule. La paire \((R,q)\) déploie l'histoire. Sous composition, la retenue obéit à une loi de cocycle ; elle n'est donc pas une annotation auxiliaire, mais la donnée qui permet de recomposer une longue opération à partir de ses étapes.

\begin{example}[Même résidu, histoire différente]
Dans la carte nonadique,
\[
 8+8=7+9\cdot1,\qquad 8\cdot2=7+9\cdot1.
\]
Les deux opérations publient le même résidu et le même quotient local, mais diffèrent par le feuillet. L'étiquette \(\Sigma/\Pi\) reste nécessaire. Dans les voies plus longues, l'ordre de ces feuillets modifie la retenue composée et le bord atteint.
\end{example}

\section{Incompatibilité opératorielle des feuillets additif et multiplicatif}

Pour \(d\ge2\), soit \(\mathscr A_d=E_9^d\). Les espérances conditionnelles \(P_{\Sigma,d}\) et \(P_{\Pi,d}\) sur les deux évaluations satisfont
\[
 \|[P_{\Sigma,2},P_{\Pi,2}]\|=\frac{\sqrt2}{3},
 \qquad
 \|[P_{\Sigma,3},P_{\Pi,3}]\|=\frac{\sqrt3}{4}.
\]
En dimension trois, le mode tritique de phase opératoire possède la valeur singulière \(s=1/2\), et
\[
 s^2=\frac14,\qquad 1-s^2=\frac34=\frac{90}{120},
 \qquad s\sqrt{1-s^2}=\frac{\sqrt3}{4}.
\]
Ainsi, la séparation entre les feuillets n'est pas ajoutée après coup : elle est quantifiée dans le support APP lui-même.

\section{Courbure discrète somme--produit}

En résidus, écrivons
\[
 S(x,y)=x+y,\qquad P(x,y)=xy\pmod9.
\]
Pour un potentiel \(T\), on définit les différences progressives
\[
 \Delta_xT(x,y)=T(x+1,y)-T(x,y),\qquad
 \Delta_yT(x,y)=T(x,y+1)-T(x,y).
\]
Avec \(A_x^T=\Delta_xT\) et \(A_y^T=\Delta_yT\), la courbure d'un mélange ordonné est
\[
 F(T_x,T_y)=\Delta_xA_y^{T_y}-\Delta_yA_x^{T_x}.
\]
Comme
\[
 \Delta_xS=\Delta_yS=1,\qquad
 \Delta_xP=y,\qquad \Delta_yP=x,
\]
on obtient exactement
\begin{equation}
 F(S,S)=F(P,P)=0,\qquad
 F(S,P)=+1,\qquad F(P,S)=-1.
 \label{espiral:eq:curvatura-mixta-app}
\end{equation}

\begin{theorem}[Courbure mixte orientée d'APP]
Les compositions pures des feuillets d'APP sont plates dans la convention de plaquettes précédente ; les compositions mixtes possèdent une courbure constante de signe opposé selon l'ordre.
\end{theorem}
\begin{proof}
En substituant les quatre différences dans la définition,
\[
 \Delta_x\Delta_yS-\Delta_y\Delta_xS=0,
 \quad
 \Delta_x\Delta_yP-\Delta_y\Delta_xP=0
\]
pour les deux compositions pures, tandis que
\[
 \Delta_x\Delta_yP-\Delta_y\Delta_xS=1,
 \qquad
 \Delta_x\Delta_yS-\Delta_y\Delta_xP=-1.
\]
L'identité se vérifie dans chacune des \(81\) plaquettes.
\end{proof}

Pour un rectangle de coordonnées de côtés \(a,b\), l'annulation des arêtes intérieures donne la forme finie de Stokes :
\[
 W_{\partial R_{a,b}}(S,P)=ab,\qquad
 W_{\partial R_{a,b}}(P,S)=-ab\pmod9.
\]

\begin{figure}[htbp]
\centering
\begin{tikzpicture}[scale=.72]
\foreach \x in {0,...,8}{
  \draw[HMTLine] (\x,0)--(\x,8);
  \draw[HMTLine] (0,\x)--(8,\x);
}
\draw[HMTNavy,thick] (0,0) rectangle (8,8);
\draw[->,HMTBlue,very thick] (2,2)--(5,2) node[midway,below] {$S$};
\draw[->,HMTRed,very thick] (5,2)--(5,5) node[midway,right] {$P$};
\draw[->,HMTBlue!70!black,very thick] (5,5)--(2,5);
\draw[->,HMTRed!80!black,very thick] (2,5)--(2,2);
\node[fill=white,inner sep=2pt] at (3.5,3.5) {$F(S,P)=+1$};
\node[below,HMTGray,font=\small\sffamily] at (4,-.75) {soporte APP \(9\times9\)};
\end{tikzpicture}
\caption{La courbure naît de l'ordre des feuillets sur un même support.}
\label{espiral:fig:curvatura-app}
\end{figure}

\section{Antécédents APP de la classification radiale}

APP apporte trois ingrédients distincts :
\begin{enumerate}
\item un domaine bidimensionnel antérieur à \(\R^2\) ;
\item deux opérations dont le mélange ordonné distingue l'orientation ;
\item un registre de résidu, quotient, feuillet et retenue capable de mémoriser le chemin qui a produit une même sortie.
\end{enumerate}
La courbure \(\pm1\) est l'antécédent local de l'orientation radiale, mais elle ne doit pas être identifiée sans preuve au degré global \(p\). Le passage de l'une à l'autre requiert l'accumulation TPK et le cocycle déclaré dans la \cref{espiral:chap:cociclos-radiales}.

\chapter{TRIT : orientation locale et 3 régimes géométriques}

\section{État ternaire relevé}

L'observable visible est
\[
 \TRIT=\{-1,0,+1\}.
\]
Dans une étape locale avec une retenue, la décomposition équilibrée s'écrit
\[
 \widehat\tau=(r,c),\qquad a+b=r+3c,
 \qquad r\in\{-1,0,+1\},\ c\in\Z.
\]
Le zéro visible possède au moins trois relèvements élémentaires :
\[
 0^+=(0,+1),\qquad0^0=(0,0),\qquad0^-=(0,-1).
\]
L'égalité visible ne fixe pas la continuation. Celle-ci dépend aussi du feuillet, de l'orientation, de la hauteur, de la retenue, du bord, de la voie et de la phase.

\section{Algèbres quadratiques}

Para \(\tau\in\TRIT\), sea
\[
 \mathbb A_\tau=\R[J]/(J^2+\tau I).
\]
On obtient trois régimes :
\[
 J_{+1}^2=-I,\qquad J_0^2=0,\qquad J_{-1}^2=+I,
\]
et une seule identité exponentielle
\[
 \exp(tJ_\tau)=C_\tau(t)I+S_\tau(t)J_\tau,
\]
où
\[
\begin{array}{c|c|c|c}
\tau & J_\tau^2 & C_\tau(t)&S_\tau(t)\\ \hline
+1&-I&\cos t&\sin t\\
0&0&1&t\\
-1&I&\cosh t&\sinh t
\end{array}
\]

\begin{figure}[htbp]
\centering
\begin{tikzpicture}[>=Latex]
\begin{scope}[xshift=-5cm]
 \draw[->,HMTGray] (-1.8,0)--(1.8,0);
 \draw[->,HMTGray] (0,-1.8)--(0,1.8);
 \draw[HMTBlue,very thick] (0,0) circle (1.25);
 \draw[->,HMTRed,thick] (0,0)--({1.25*cos(40)},{1.25*sin(40)});
 \node[below,HMTNavy,font=\sffamily\bfseries] at (0,-2.05) {elliptique \(J^2=-I\)};
\end{scope}
\begin{scope}
 \draw[->,HMTGray] (-1.8,0)--(1.8,0);
 \draw[->,HMTGray] (0,-1.8)--(0,1.8);
 \draw[HMTBlue,very thick] (-1.3,-.65)--(1.3,.65);
 \draw[->,HMTRed,thick] (0,0)--(1.1,.55);
 \node[below,HMTNavy,font=\sffamily\bfseries] at (0,-2.05) {parabolique \(J^2=0\)};
\end{scope}
\begin{scope}[xshift=5cm]
 \draw[->,HMTGray] (-1.8,0)--(1.8,0);
 \draw[->,HMTGray] (0,-1.8)--(0,1.8);
 \draw[HMTBlue,very thick,domain=-1.3:1.3,samples=80] plot ({cosh(\x)},{sinh(\x)});
 \draw[HMTBlue,very thick,domain=-1.3:1.3,samples=80] plot ({-cosh(\x)},{sinh(\x)});
 \draw[->,HMTRed,thick] (0,0)--(1.25,.75);
 \node[below,HMTNavy,font=\sffamily\bfseries] at (0,-2.05) {hyperbolique \(J^2=I\)};
\end{scope}
\end{tikzpicture}
\caption{Les trois régimes quadratiques du TRIT. Ce sont des types de transport, non des noms de courbes concrètes.}
\label{espiral:fig:trit-regimenes}
\end{figure}

\section{Séparation entre régime TRIT, caractère radial et courbure de Frenet}

La théorie requiert trois niveaux qui ne doivent pas être aplatis :
\begin{enumerate}
\item \(\tau\) type le régime local de transport ;
\item \(p\) type le caractère radial accumulé d'une histoire ;
\item \(\kappa_{\mathrm{F}}\) et \(\mathcal T_{\mathrm{F}}\) décrivent la courbe déjà publiée dans un espace réel.
\end{enumerate}
Un cercle possède une holonomie angulaire elliptique, mais le présent TRIT \(0\) est un seuil parabolique. Le cercle et le présent ne sont pas le même objet. De même, une spirale logarithmique publiée peut inclure des segments dont les transitions locales parcourent les trois types TRIT.

\begin{principle}[Séparation des niveaux]
Aucune identité entre \(\tau\), \(p\) et \(\kappa_{\mathrm{F}}\) n'est admise par ressemblance de signe, de nom ou de dessin. Toute relation doit être construite au moyen d'une application dont le domaine, le codomaine et l'action sont déclarés.
\end{principle}

\section{Involution d'orientation}

Pour un mot \(\boldsymbol\tau=(\tau_1,\ldots,\tau_n)\),
\[
 C_{\mathrm{rev}}(\boldsymbol\tau)
 =(-\tau_n,\ldots,-\tau_1),
 \qquad C_{\mathrm{rev}}^2=\id.
\]
L'involution échange les orientations et préserve le seuil. Dans le secteur des voies réversibles, elle induira l'inversion de l'holonomie affine, mais cette conclusion exige que la représentation locale satisfasse \(h_{e^\dagger}=h_e^{-1}\).

\chapter{État enrichi, continu unique et double projection}

\section{Système inverse d'histoires}

Soient \(\mathsf S_n\) les états admissibles à la profondeur \(n\) et
\[
 \pi_{n+1,n}:\mathsf S_{n+1}\longrightarrow\mathsf S_n
\]
les troncatures. L'espace des histoires complètes est
\[
 X_\infty=\varprojlim_n\mathsf S_n.
\]
Avec l'ultramétrique
\[
 d(x,y)=3^{-m(x,y)},
\]
où \(m(x,y)\) est la première profondeur à laquelle elles diffèrent, les cylindres à préfixe fixé sont ouverts et fermés. La structure du continu s'obtient en construisant conjointement, à partir du même état APP--TRIT--TPK, ses cinq constructions consubstantielles, qui émergent corrélativement, et en passant par la terminalité, la naturalité et la limite. Cette publication ne les dissocie pas et ne les redéfinit pas au moyen de courbes.

\section{Évaluation archimédienne au moyen de cylindres compatibles}

Soit \(\widehat{\mathbb Q}\) la complétion uniforme de \(\mathbb Q\), construite comme classes de suites rationnelles de Cauchy. La construction calibrée associe à chaque état fini une adresse entière \(k_n(s)\) et le cylindre rationnel semi-fermé
\[
 I_n(s)=\{q\in\mathbb Q:k_n(s)3^{-n}\le q<(k_n(s)+1)3^{-n}\}.
\]
Il satisfait simultanément :
\begin{enumerate}[label=\textup{(A\arabic*)}]
\item \(\bigcup_{s\in\mathsf S_n}I_n(s)=\mathbb Q\) pour tout \(n\) ;
\item si \(t\in\mathsf S_{n+1}\) se tronque en \(s\), alors \(I_{n+1}(t)\subseteq I_n(s)\), et les enfants partitionnent le parent selon la convention semi-fermée de bord ;
\item il existe une constante \(C>0\) et \(0<\lambda<1\) telles que
\[
 \operatorname{diam}I_n(s)\le C\lambda^n,
 \qquad 0<\lambda<1;
\]
\item toute chaîne de cylindres emboîtés provient d'une histoire compatible, avec les deux expansions de bord identifiées par la relation archimédienne.
\end{enumerate}

\begin{theorem}[Fermeture archimédienne du continu généalogique]
\label{thm:espiral-cierre-arquimediano-continuo-genealogico}
Chaque histoire calibrée détermine une unique valeur
\[
 \{\operatorname{val}(x)\}
 =\bigcap_n\overline{I_n(x_n)}^{\widehat{\mathbb Q}}.
\]
L'évaluation
\[
 \operatorname{val}:X_\infty^{\mathrm{cal}}\longrightarrow\widehat{\mathbb Q}
\]
est continue, höldérienne et surjective. Le quotient par égalité de valeur satisfait
\[
 X_\infty^{\mathrm{cal}}/{\sim_{\arq}}
 \cong\widehat{\mathbb Q}\cong\R,
\]
où la dernière identification est la reconnaissance archimédienne usuelle de la complétion de \(\mathbb Q\), non une entrée du générateur.
\end{theorem}
\begin{proof}
La contraction des diamètres et la complétude de \(\widehat{\mathbb Q}\) réduisent chaque chaîne emboîtée à un point. La compatibilité des cylindres démontre la continuité ; la borne géométrique des diamètres donne l'estimation de Hölder. Par (A1) et (A2), chaque classe de Cauchy rationnelle sélectionne à toute profondeur un cylindre qui la contient ; (A4) fournit une histoire compatible et résout les bords de manière unique. Il en résulte la surjectivité. Deux histoires ont la même image exactement lorsqu'elles appartiennent à la relation \(\sim_{\arq}\), si bien que l'application induite sur le quotient est une bijection continue et ouverte sur la topologie finale des cylindres ; c'est donc un homéomorphisme.
\end{proof}

La droite réelle est donc une sortie démontrée de l'évaluation et non un domaine présupposé. L'histoire ne se réduit pas à sa valeur.

\section{Double projection}

À partir d'un même domaine calibré, on distingue
\[
 \mathfrak R:X_\infty^{\mathrm{cal}}\longrightarrow\R,
 \qquad
 \mathfrak E:X_\infty^{\mathrm{cal}}\times\mathscr C_{\mathrm{exc}}
 \longrightarrow\mathbf{Exc}.
\]
La première contracte les cylindres compatibles ; la seconde conserve l'incidence de bord vers le corridor exceptionnel. La théorie spirale s'insère dans la première branche, mais conserve dans son registre les données nécessaires pour comparer la seconde. On n'identifie pas la valeur à l'incidence et l'on ne fait pas générer rétrospectivement le continu par une publication.

\begin{figure}[htbp]
\centering
\begin{tikzpicture}[>=Latex,node distance=12mm and 13mm,
 n/.style={draw=HMTNavy,fill=white,minimum height=9mm,text width=34mm,align=center,font=\small\sffamily}]
\node[n] (x) {$X_\infty^{\mathrm{cal}}$\\état enrichi commun};
\node[n,below left=of x] (r) {publication archimédienne\\valeur et courbe dans \(\R^d\)};
\node[n,below right=of x] (e) {publication exceptionnelle\\incidence et automorphismes};
\draw[->,HMTBlue,very thick] (x)--node[left,font=\small] {$\mathfrak R$} (r);
\draw[->,HMTRed,very thick] (x)--node[right,font=\small] {$\mathfrak E_c$} (e);
\draw[HMTGold,dashed,thick] (r)--node[below,font=\scriptsize\sffamily]{invariantes comunes} (e);
\end{tikzpicture}
\caption{Les deux projections sont corrélatives et postérieures au même état.}
\label{espiral:fig:doble-proyeccion-espiral}
\end{figure}

\section{5 régions de \texorpdfstring{\(\pi\)}{pi} et multiplicité des généalogies}

La microfibre de \(\pi\) contient cinq régions distinctes de multiplicités
\[
 432+4\cdot144=1008.
\]
La même valeur visible conserve des histoires régionales différentes. Ce fait fournit le précédent exact de la classification spirale : l'égalité d'une courbe publiée n'impose pas l'égalité du feuillet, de la retenue, du bord ni de la mémoire. La théorie n'invente pas cette multiplicité ; elle la transporte vers un nouveau type de sortie géométrique.

\begin{resultbox}[title={Antériorité causale du continu discret par rapport au lecteur radial}]
Jusqu'ici, on a construit le domaine, les feuillets, les régimes, l'holonomie, l'état enrichi et la publication du continu. Ce n'est que maintenant qu'il est licite d'introduire des coordonnées réelles \((r,\theta)\) et de reconnaître des courbes conventionnelles.
\end{resultbox}
\part{Théorie généalogique des correspondances spirales}

\chapter{Enveloppe affine universelle des feuillets APP}

\section{Construction du module radial universel et de son enveloppe endomorphique}

Les opérations \(u\mapsto u+a\) et \(u\mapsto\lambda u\) suggèrent le groupe affine de la droite. Cependant, APP contient des secteurs multiplicatifs non inversibles et précède la publication réelle. La cible interne appropriée est donc le monoïde des endomorphismes affines d'un module radial universel.

Soit \(\widetilde X_{\APP}\) l'ensemble des états APP complets, c'est-à-dire des états qui conservent les deux feuillets, le résidu, le quotient et la retenue, et ajoutons un symbole absorbant \(\bot\). Nous définissons le module libre
\begin{equation}
 U_{\APP}:=\Z^{(\widetilde X_{\APP}\sqcup\{\bot\})}/\langle[\bot]\rangle.
 \label{espiral:eq:modulo-radial-universal}
\end{equation}
Pour chaque transition enrichie \(e\), de source \(x_e\) et de chiffres actifs \((i_e,j_e)\), APP produit sans évaluation décimale les enregistrements
\[
 i_e+j_e=R_e^++9q_e^+,
 \qquad
 i_ej_e=R_e^\times+9q_e^\times.
\]
Soient \(s_e^\Sigma(x_e)\) et \(s_e^\Pi(x_e)\) les états complets obtenus en adjoignant, respectivement, \((R_e^+,q_e^+)\) et \((R_e^\times,q_e^\times)\), en conservant l'autre feuillet dans le registre. Le relèvement additif détermine le vecteur explicite
\[
 A_e=[s_e^\Sigma(x_e)]-[x_e]\in U_{\APP},
\]
et le relèvement multiplicatif détermine l'application totale \(m_e:\widetilde X_{\APP}\sqcup\{\bot\}\to
\widetilde X_{\APP}\sqcup\{\bot\}\), prolongée par
\[
 m_e(x)=
 \begin{cases}
 s_e^\Pi(x_e),&x=x_e,\\
 \bot,&x\neq x_e,
 \end{cases}
 \qquad m_e(\bot)=\bot.
\]
Sa linéarisation est
\[
 M_e([x])=[m_e(x)]\in U_{\APP}.
\]
Ainsi, \(A_e\) et \(M_e\) ne sont pas des paramètres libres : ce sont les deux publications linéarisées de la même transition APP complète.

On définit
\[
 \AffEnd(U_{\APP})
 =\operatorname{End}(U_{\APP})\ltimes U_{\APP}
\]
avec la composition
\begin{equation}
 (\Lambda_2,C_2)\circ(\Lambda_1,C_1)
 =\bigl(\Lambda_2\Lambda_1,
 C_2+\Lambda_2C_1\bigr).
 \label{espiral:eq:producto-affend}
\end{equation}
La paire agit par \(u\mapsto\Lambda u+C\). La partie linéaire peut être non inversible. Le sous-groupoïde des unités s'obtient en restreignant \(\Lambda\in\operatorname{Aut}(U_{\APP})\).

\begin{proposition}[Propriété universelle]
Toute application de l'ensemble des états complets dans un module \(V\), nulle en \(\bot\), se prolonge de manière unique en un morphisme \(\chi:U_{\APP}\to V\). Si elle entrelace en outre les relèvements multiplicatifs, \(\chi M_e=M_e^V\chi\), elle transporte alors de manière unique les facteurs affines \(h_e\) et toutes leurs holonomies préfixes.
\end{proposition}
\begin{proof}
La première affirmation est la propriété universelle du module libre quotienté par \([\bot]=0\). La seconde découle de l'application de \(\chi\) à la composition semi-directe : \(\chi(C_2+\Lambda_2C_1)=
\chi(C_2)+\Lambda_2^V\chi(C_1)\).
\end{proof}

\section{Représentation locale d'une transition}

Soit \(e:x\to y\) une transition TPK. Le sélecteur tritique de phase opératoire détermine si le feuillet additif, le feuillet multiplicatif ou le seuil agit. On définit
\begin{equation}
 h_e=
 \begin{cases}
 (I,A_e),&\text{feuillet }\Sigma,\\
 (M_e,0),&\text{feuillet }\Pi,\\
 (I,0),&\text{seuil},
 \end{cases}
 \label{espiral:eq:factor-afin-local}
\end{equation}
où \(A_e\) et \(M_e\) sont les relèvements APP complets construits ci-dessus. On ne prend pas seulement leurs résidus : le quotient et la retenue restent dans le registre.

\begin{definition}[Représentation radiale interne]
Une représentation radiale interne est une affectation
\[
 \rho_{\rad}:\mathscr P_{\TPK}^{\enr}
 \longrightarrow\AffEnd(U_{\APP})
\]
qui envoie chaque flèche sur \(h_e\), l'identité sur \((I,0)\) et la composition des chemins sur le produit \cref{espiral:eq:producto-affend}.
\end{definition}

La définition ne consulte aucune courbe. Les facteurs proviennent du feuillet et du relèvement que l'état contient déjà. Une réalisation numérique ultérieure peut envoyer \(U_{\APP}\) sur \(\R\), mais cette évaluation ne modifie pas la représentation interne.

\section{Holonomies préfixes}

Pour
\[
 \gamma_N=e_N\cdots e_1,
\]
nous définissons
\[
 H_0=(I,0),\qquad
 H_j=h_{e_j}\cdots h_{e_1}=(\Lambda_j,C_j).
\]

\begin{theorem}[Holonomie affine fonctorielle]
Pour des voies composables \(\gamma_1,\gamma_2\),
\[
 H_{\gamma_2\gamma_1}=H_{\gamma_2}H_{\gamma_1}.
\]
Si \(H_{\gamma_i}=(\Lambda_i,C_i)\), alors
\begin{equation}
 \Lambda_{\gamma_2\gamma_1}=\Lambda_2\Lambda_1,
 \qquad
 C_{\gamma_2\gamma_1}=C_2+\Lambda_2C_1.
 \label{espiral:eq:holonomia-afin-functorial}
\end{equation}
\end{theorem}
\begin{proof}
L'action successive sur \(u\) est
\[
 u\xmapsto{H_{\gamma_1}}\Lambda_1u+C_1
 \xmapsto{H_{\gamma_2}}
 \Lambda_2\Lambda_1u+C_2+\Lambda_2C_1.
\]
L'unicité de la décomposition affine donne \cref{espiral:eq:holonomia-afin-functorial}.
\end{proof}

Pour un bloc de neuf étapes,
\begin{align}
 \Lambda_9&=\Lambda_{e_9}\cdots\Lambda_{e_1},\\
 C_9&=C_{e_9}+\Lambda_{e_9}C_{e_8}
 +\Lambda_{e_9}\Lambda_{e_8}C_{e_7}+\cdots
 +\Lambda_{e_9}\cdots\Lambda_{e_2}C_{e_1}.
 \label{espiral:eq:traslacion-nueve}
\end{align}
La position de chaque facteur importe. C'est la forme affine de la mémoire d'ordre.

\section{Commutateur somme--produit}

Dans le secteur inversible et sur une réalisation unidimensionnelle, soient
\[
 A_a(u)=u+a,\qquad M_\lambda(u)=\lambda u.
\]
Alors
\begin{equation}
 A_aM_\lambda A_a^{-1}M_\lambda^{-1}
 =A_{a(1-\lambda)}.
 \label{espiral:eq:conmutador-afin}
\end{equation}
La preuve est directe :
\[
 u\mapsto\lambda^{-1}u\mapsto\lambda^{-1}u-a
 \mapsto u-\lambda a\mapsto u+a(1-\lambda).
\]
L'identité n'est affirmée que lorsque \(\lambda\) est une unité. Dans le monoïde complet, le défaut de commutation est conservé au moyen des facteurs et du registre, non au moyen d'un commutateur qui requiert des inverses inexistants.

\begin{figure}[!htb]
\centering
\begin{tikzpicture}[>=Latex,node distance=18mm and 28mm,
 n/.style={draw=HMTNavy,fill=white,minimum width=23mm,minimum height=8mm,font=\small}]
\node[n] (u) {$u$};
\node[n,right=of u,yshift=10mm] (a) {$u+a$};
\node[n,right=of a] (ma) {$\lambda u+\lambda a$};
\node[n,right=of u,yshift=-10mm] (m) {$\lambda u$};
\node[n,right=of m] (am) {$\lambda u+a$};
\draw[->,HMTBlue,thick] (u)--node[above left] {$A_a$} (a);
\draw[->,HMTRed,thick] (a)--node[above] {$M_\lambda$} (ma);
\draw[->,HMTRed,thick] (u)--node[below left] {$M_\lambda$} (m);
\draw[->,HMTBlue,thick] (m)--node[below] {$A_a$} (am);
\draw[->,HMTGold,dashed,thick] (ma)--node[right,font=\scriptsize]
  {$A_{a(1-\lambda)}$} (am);
\node[below=12mm of u,HMTGray,font=\small\sffamily,text width=80mm,align=center]
{Les deux ordres aboutissent à des points distincts ; la flèche verticale mesure exactement le défaut de composition.};
\end{tikzpicture}
\caption{La différence somme--produit se manifeste comme une holonomie affine.}
\label{espiral:fig:cuadrado-afin}
\end{figure}

\chapter{Cocycles radiaux dans le relèvement TRIT de 2 demi-couronnes}
\label{espiral:chap:cociclos-radiales}

\section{Catégorie relevée et caractères élémentaires}

Soit \(\mathscr P_{\TPK}^{2\mathrm c,\enr}\) la catégorie de chemins dont les flèches sont des triplets
\[
 \widehat e=(e,\tau_+(e),\tau_-(e)),
 \qquad \tau_\pm(e)\in\{-1,0,+1\},
\]
où la paire orientée appartient à l'état enrichi et le foncteur d'oubli envoie \(\widehat e\) sur la transition TPK \(e\). Nous définissons sur chaque générateur
\begin{equation}
 p(\widehat e)=\tau_+(e)+\tau_-(e),
 \qquad
 a(\widehat e)=\tau_+(e)-\tau_-(e).
 \label{espiral:eq:carta-doble-trit}
\end{equation}
Pour une voie \(\widehat\gamma_N=\widehat e_N\cdots\widehat e_1\), les cocycles accumulés sont
\begin{equation}
 P_j=\sum_{k=1}^{j}p(\widehat e_k),
 \qquad
 A_j=\sum_{k=1}^{j}a(\widehat e_k),
 \qquad P_0=A_0=0.
 \label{espiral:eq:cociclos-radiales}
\end{equation}
L'additivité sous concaténation démontre que \((P,A)\) est un cocycle à valeurs dans \(\Z^2\) de la catégorie relevée. Cette affirmation n'identifie pas par décret la paire à un autre cocycle précédemment nommé : elle la construit à partir des deux orientations effectivement conservées dans chaque flèche. La subdivision par microflèches immobiles préserve les deux sommes.

\section{Carte locale exhaustive des caractères radial et orienté}

\begin{longtable}{cccccc}
\caption{Carte locale complète des deux orientations TRIT.}
\label{espiral:tab:doble-trit}\\
\toprule
\(\tau_+\)&\(\tau_-\)&\(p\)&\(a\)&\(r_3(p)\)&\(c_3(p)\)\\
\midrule
\endfirsthead
\toprule
\(\tau_+\)&\(\tau_-\)&\(p\)&\(a\)&\(r_3(p)\)&\(c_3(p)\)\\
\midrule
\endhead
+1&+1&+2&0&-1&1\\
+1&0&+1&+1&+1&0\\
+1&-1&0&+2&0&0\\
0&+1&+1&-1&+1&0\\
0&0&0&0&0&0\\
0&-1&-1&+1&-1&0\\
-1&+1&0&-2&0&0\\
-1&0&-1&-1&-1&0\\
-1&-1&-2&0&+1&-1\\
\bottomrule
\end{longtable}

Les trois états avec \(p=0\) se distinguent par \(a=2,0,-2\). Le degré radial ne suffit donc pas à reconstruire l'orientation.

\section{Décomposition équilibrée}

Pour tout \(p\in\Z\), il existe une unique décomposition
\[
 p=r_3(p)+3c_3(p),
 \qquad r_3(p)\in\{-1,0,+1\},\quad c_3(p)\in\Z.
\]
Dans le bloc local \(p\in[-2,2]\), la dernière colonne de \cref{espiral:tab:doble-trit} donne la retenue. Dans les voies longues, \(c_3(p)\) n'est pas restreint à \(\{-1,0,+1\}\). Cette précision évite de transformer un exemple local en une fausse loi globale.

\section{Involution}

L'involution locale
\[
 (\tau_+,\tau_-)\longmapsto(-\tau_-,-\tau_+)
\]
produit
\[
 p\longmapsto-p,
 \qquad
 a\longmapsto a.
\]
Le caractère radial change d'orientation, tandis que la différence orientée demeure. Cette séparation sera visible dans l'identité \(L_{-p}(x^{-1})=-L_p(x)\).

\chapter{Coordonnée puissance--logarithmique et loi somme--produit}

\section{Définition et domaine}

Pour \(x>0\), nous définissons
\begin{equation}
 L_p(x)=
 \begin{cases}
 \dfrac{x^p-1}{p},&p\neq0,\\[2mm]
 \log x,&p=0.
 \end{cases}
 \label{espiral:eq:lp-def}
\end{equation}
Son inverse est
\begin{equation}
 \exp_p(u)=
 \begin{cases}
 (1+pu)^{1/p},&p\neq0,\quad1+pu>0,\\[1mm]
 \e^u,&p=0.
 \end{cases}
 \label{espiral:eq:exp-p}
\end{equation}
La condition \(1+pu>0\) n'est pas ornementale : elle fixe la branche réelle positive.

\section{Addition déformée}

Dans le domaine correspondant, soit
\[
 u\boxplus_pv=u+v+puv.
\]

\begin{theorem}[Loi somme--produit]
Pour \(x,y>0\),
\begin{equation}
 L_p(xy)=L_p(x)\boxplus_pL_p(y).
 \label{espiral:eq:lp-producto}
\end{equation}
De plus, \(0\) est l'identité de \(\boxplus_p\), et
\[
 \exp_p(u\boxplus_pv)=\exp_p(u)\exp_p(v).
\]
\end{theorem}
\begin{proof}
Para \(p\neq0\),
\[
 1+pL_p(xy)=(xy)^p=x^py^p
 =(1+pL_p(x))(1+pL_p(y)).
\]
En développant le dernier produit et en divisant par \(p\), on obtient \cref{espiral:eq:lp-producto}. Pour \(p=0\), c'est l'identité \(\log(xy)=\log x+\log y\). L'égalité pour \(\exp_p\) découle de l'application de l'inverse.
\end{proof}

Cette loi exprime exactement l'articulation APP : le produit dans la coordonnée radiale est publié comme une somme déformée dans la coordonnée \(L_p\). Le terme \(puv\) conserve l'interaction qui disparaîtrait dans une somme ordinaire.

\section{Inversion bidirectionnelle}

\begin{proposition}[Identité d'inversion]
Pour \(p\in\Z\) et \(x>0\),
\begin{equation}
 L_{-p}(x^{-1})=-L_p(x).
 \label{espiral:eq:lp-inversion}
\end{equation}
\end{proposition}
\begin{proof}
Si \(p\neq0\),
\[
 L_{-p}(x^{-1})
 =\frac{(x^{-1})^{-p}-1}{-p}
 =-\frac{x^p-1}{p}.
\]
Pour \(p=0\), \(\log(x^{-1})=-\log x\).
\end{proof}

L'involution échange simultanément expansion/contraction et \(p/-p\). Elle n'élimine pas la mémoire d'orientation, qui demeure dans \(a\).

\section{Limite logarithmique}

Para \(x>0\),
\[
 \lim_{p\to0}\frac{x^p-1}{p}=\log x.
\]
En effet, \(x^p=\e^{p\log x}=1+p\log x+O(p^2)\). La branche logarithmique n'est donc pas une exception ajoutée : c'est la limite continue de la famille. Dans la théorie généalogique, cependant, \(p=0\) continue de conserver \(a\) et le registre ; la limite fonctionnelle n'identifie pas ses trois origines locales de \cref{espiral:tab:doble-trit}.

\chapter{Classification puissance--logarithmique des courbes publiées}

\section{Translation uniforme dans la coordonnée radiale}

Supposons, après la publication du continu, que
\[
 L_p(r/r_0)=\beta m,
 \qquad
 \theta=\theta_0+\omega m,
 \qquad \omega\neq0.
\]
Eliminando \(m\),
\begin{equation}
 r_p(\theta)=r_0
 \left[1+p\frac{\beta}{\omega}
 (\theta-\theta_0)\right]^{1/p},
 \qquad p\neq0,
 \label{espiral:eq:familia-potencia}
\end{equation}
et
\begin{equation}
 r_0(\theta)=r_0
 \exp\!\left[\frac{\beta}{\omega}
 (\theta-\theta_0)\right].
 \label{espiral:eq:familia-log}
\end{equation}

\begin{theorem}[Publication des cinq classes canoniques]
Après changement d'origine de \(\theta\) et changement d'échelle de \(r\), les spécialisations de \cref{espiral:eq:familia-potencia,espiral:eq:familia-log} sont :
\[
\begin{array}{c|c|c}
p&\text{équation normalisée}&\text{reconnaissance conventionnelle}\\ \hline
2&r^2=a^2\theta&\text{spirale de Fermat}\\
1&r=a\theta&\text{spirale d'Archimède}\\
0&r=r_0\e^{b\theta}&\text{spirale logarithmique}\\
-1&r=a/\theta&\text{spirale hyperbolique}\\
-2&r^2=a^2/\theta&\text{lituus}
\end{array}
\]
dans les domaines réels positifs correspondants.
\end{theorem}
\begin{proof}
Pour \(p=1\), \(r\) est affine en \(\theta\), et le changement d'origine élimine le terme constant. Pour \(p=2\), \(r^2\) est affine. Pour \(p=-1\), \(r^{-1}\) est affine ; pour \(p=-2\), \(r^{-2}\) est affine. La branche \(p=0\) est \cref{espiral:eq:familia-log}. Les constantes positives sont absorbées dans \(a,b\) sans altérer l'exposant.
\end{proof}

\begin{figure}[!htb]
\centering
\begin{tikzpicture}[scale=.62]
\begin{scope}[xshift=0cm,yshift=0cm]
  \draw[->,HMTLine] (-1.9,0)--(1.9,0);
  \draw[->,HMTLine] (0,-1.9)--(0,1.9);
  \draw[HMTBlue,very thick,domain=.15:18.2,samples=180,smooth,variable=\t]
    plot ({sqrt(.14*(\t+1))*cos(deg(\t))},{sqrt(.14*(\t+1))*sin(deg(\t))});
  \node[below,text=HMTGray,font=\scriptsize\sffamily] at (0,-2.1) {Fermat, \(p=2\)};
\end{scope}
\begin{scope}[xshift=5cm,yshift=0cm]
  \draw[->,HMTLine] (-1.9,0)--(1.9,0);
  \draw[->,HMTLine] (0,-1.9)--(0,1.9);
  \draw[HMTNavy,very thick,domain=.15:18.2,samples=180,smooth,variable=\t]
    plot ({.08*(\t+1)*cos(deg(\t))},{.08*(\t+1)*sin(deg(\t))});
  \node[below,text=HMTGray,font=\scriptsize\sffamily] at (0,-2.1) {Archimède, \(p=1\)};
\end{scope}
\begin{scope}[xshift=10cm,yshift=0cm]
  \draw[->,HMTLine] (-1.9,0)--(1.9,0);
  \draw[->,HMTLine] (0,-1.9)--(0,1.9);
  \draw[HMTRed,very thick,domain=.15:18.2,samples=180,smooth,variable=\t]
    plot ({.18*exp(.055*\t)*cos(deg(\t))},{.18*exp(.055*\t)*sin(deg(\t))});
  \node[below,text=HMTGray,font=\scriptsize\sffamily] at (0,-2.1) {Logarithmique, \(p=0\)};
\end{scope}
\begin{scope}[xshift=2.5cm,yshift=-5.2cm]
  \draw[->,HMTLine] (-1.9,0)--(1.9,0);
  \draw[->,HMTLine] (0,-1.9)--(0,1.9);
  \draw[HMTGold,very thick,domain=.15:18.2,samples=180,smooth,variable=\t]
    plot ({1.8/(\t+1.8)*cos(deg(\t))},{1.8/(\t+1.8)*sin(deg(\t))});
  \node[below,text=HMTGray,font=\scriptsize\sffamily] at (0,-2.1) {Hyperbolique, \(p=-1\)};
\end{scope}
\begin{scope}[xshift=7.5cm,yshift=-5.2cm]
  \draw[->,HMTLine] (-1.9,0)--(1.9,0);
  \draw[->,HMTLine] (0,-1.9)--(0,1.9);
  \draw[HMTGray,very thick,domain=.15:18.2,samples=180,smooth,variable=\t]
    plot ({sqrt(2.0/(\t+1.8))*cos(deg(\t))},{sqrt(2.0/(\t+1.8))*sin(deg(\t))});
  \node[below,text=HMTGray,font=\scriptsize\sffamily] at (0,-2.1) {Lituus, \(p=-2\)};
\end{scope}
\end{tikzpicture}
\caption{Cinq publications de la famille puissance--logarithmique sur des axes et échelles communs.}
\label{espiral:fig:cinco-espirales}
\end{figure}

\section{Non-exhaustivité et portée exacte}

Le théorème classe la famille engendrée par translation uniforme dans la coordonnée \(L_p\). Il ne classe pas toutes les courbes planes possibles. La taxonomie généalogique complète doit inclure au moins
\[
 (p,a,[\Lambda,C],[\kappa_9],
 \omega_{\mathrm{ang}},\mu_{\mem},\varepsilon_{\mathrm{tw}}).
\]
Deux voies ayant le même \(p\) peuvent différer par l'auto-échelle, l'orientation, la retenue ou le bord et publier des courbes distinctes. Deux voies dont toutes les données visibles sont égales peuvent encore différer par une coordonnée de registre éliminée.

\section{Holonomie affine générale}

Si la coordonnée \(u=L_p(r/r_0)\) satisfait
\[
 u_{n+1}=\Lambda u_n+C,
\]
l'itération interne exacte, pour \(n\in\N\), est
\[
 u_n=\Lambda^n u_0+\sum_{k=0}^{n-1}\Lambda^kC.
\]
Si \(I-\Lambda\) est inversible, on peut écrire
\[
 u_n=u_*+\Lambda^n(u_0-u_*),
 \qquad u_*=(I-\Lambda)^{-1}C.
\]
L'interpolation continue ne s'effectue qu'après fixation d'un caractère scalaire positif \(\chi:U_{\APP}\to\R\) tel que \(\chi\Lambda=\lambda\chi\), avec \(\lambda>0\). En écrivant \(u=\chi(u)\), \(C=\chi(C)\), si \(\lambda\neq1\), on a \(u_*=C/(1-\lambda)\). Pour \(\theta_n=\theta_0+n\omega\), \(\omega\neq0\), la réalisation scalaire admet
\begin{equation}
 r(\theta)=r_0\exp_p\!\left[
 u_*+\lambda^{(\theta-\theta_0)/\omega}(u_0-u_*)
 \right].
 \label{espiral:eq:familia-afin-general}
\end{equation}
Les puissances réelles appartiennent exclusivement à cette réalisation ; la dynamique interne utilise des puissances entières de \(\Lambda\). La paire \((p,[\Lambda,C])\) élargit la liste classique : \(p\) fixe la carte radiale et la classe de conjugaison affine fixe la dynamique en son sein.

\begin{remark}[Point fixe et électron]
L'existence d'un point fixe affine n'identifie pas à elle seule ce point à l'électron central \((5,5)\). L'électron est un résultat matériel antérieur ; une identification avec un point fixe radial exige un carré de réalisation distinct.
\end{remark}

\chapter{Lecteur radial enrichi}

\section{Obstruction au lecteur réduit}

L'application initialement suggérée
\[
 X_{\TPK}^{\enr}\to(p,a,\Lambda,C,m,r,\theta)
\]
mélange les niveaux causaux et perd les préfixes. Les coordonnées \(r,\theta\) sont archimédiennes, tandis que \((\Lambda,C)\) ne conserve que le produit final.

\begin{proposition}[La paire finale n'admet pas de troncature naturelle]
Il n'existe pas en général d'application qui récupère le premier facteur d'un mot affine à partir de son produit final.
\end{proposition}
\begin{proof}
Dans la notation \((\Lambda,C)\),
\[
 (1,2)(2,0)=(2,2),
 \qquad
 (1,1)(2,1)=(2,2).
\]
Les deux produits finaux coïncident, mais leurs premiers préfixes sont \((2,0)\) et \((2,1)\). Une fonction du produit final devrait lui attribuer deux troncatures distinctes, ce qui est impossible.
\end{proof}

\section{Définition du lecteur complet}

Soit \(\widehat\gamma_N=\widehat e_N\cdots\widehat e_1\) une voie de la catégorie relevée et soit \((H_0,\Led_0)=((I,0),\Led(x_0))\) sa donnée initiale. Nous définissons
\begin{equation}
 \boxed{
 \mathcal P_{\rad,N}^{\enr}(\widehat\gamma_N)
 =\left(
 (H_0,\Led_0);
 \bigl(p(\widehat e_j),a(\widehat e_j),P_j,A_j,
 h_{e_j},H_j,\Led_j(\widehat\gamma_N)\bigr)_{j=1}^{N}
 \right).}
 \label{espiral:eq:lector-radial-enriquecido}
\end{equation}
Le codomaine \(\mathcal D_{\rad,N}^{\enr}\) contient des séquences compatibles de ces données, y compris l'arête sous-jacente stockée dans chaque entrée du registre et ses applications source et cible. La sortie conserve la donnée initiale, chaque facteur, chaque holonomie préfixe, les cocycles locaux et accumulés et chaque actualisation de mémoire.

\begin{definition}[Secteur radial homogène admissible]
Soit \(\mathcal D_{\rad,\infty}^{\mathrm{hom}}\) l'espace des données marquées
\[
 D=(D^{\enr};r_0,p,u_0,\beta,\theta_0,\omega),
\]
où \(D^{\enr}\in\mathcal D_{\rad,\infty}^{\enr}\), \(r_0>0\), \(p\in\{-2,-1,0,1,2\}\), \(u_0,\beta,\theta_0,\omega\in\R\) et \((\beta,\omega)\neq(0,0)\), telles que
\[
 p(\widehat e_j)=p,\qquad
 u(m)=u_0+\beta m,\qquad
 \theta(m)=\theta_0+\omega m.
\]
La marque fait partie de la donnée et n'est pas choisie après coup ; l'évaluation est donc une fonction bien définie.
Son intervalle réel de publication est
\[
 I_D=
 \begin{cases}
 \{m\in\R:1+p(u_0+\beta m)>0\},&p\neq0,\\
 \R,&p=0.
 \end{cases}
\]
Les histoires de degré variable restent dans le registre enrichi et sont publiées au moyen d'un atlas de cartes par morceaux ; elles ne déterminent pas une unique courbe de degré constant.
\end{definition}

\begin{definition}[Publication spirale homogène]
Après construction du continu et choix des caractères radial et angulaire compatibles, on définit
\[
 \ev_{\arq}^{\mathrm{esp}}:
 \mathcal D_{\rad,\infty}^{\mathrm{hom}}\longrightarrow
 \coprod_{D}C^\infty(I_D,\R^2)
\]
par
\[
 D\longmapsto
 \left[m\longmapsto
 \bigl(r_0\exp_{p}(u_0+\beta m)\cos(\theta_0+\omega m),
 r_0\exp_{p}(u_0+\beta m)\sin(\theta_0+\omega m)\bigr)\right].
\]
\end{definition}

Le lecteur interne et l'évaluation externe sont séparés. La première flèche conserve la généalogie ; la seconde peut en oublier une partie. La condition \((\beta,\omega)\neq(0,0)\) élimine uniquement la paramétrisation constante dégénérée ; dans \(I_D\), la courbe publiée est régulière.

\section{Conservation comme traçabilité}

La projection sur la retenue, le bord et la mémoire satisfait
\[
 \operatorname{pr}_{\mathrm{carry},frontier,memory}
 \circ\mathcal P_{\rad,N}^{\enr}
 =\Led_{{\mathrm{carry},frontier,memory},N}.
\]
Conserver ne signifie pas rendre constantes ces coordonnées. Cela signifie que :
\[
 g(K+9)=g(K),
\]
\[
 q_{\mem}(\Gamma_9x)=q_{\mem}(x)+1,
\]
\[
 B_{K+9}=\operatorname{Tra}_{\Gamma_9}(B_K),
\]
et la retenue est actualisée par \cref{espiral:eq:cociclo-carry}. Aucune actualisation n'est masquée par la courbe finale.

\chapter{Naturalité, équivariance et limite inverse}

\section{Naturalité sous troncature}

Pour \(M\le N\), soit
\[
 \pi_{N,M}:\Hist_N^{\enr}(\TPK)\to\Hist_M^{\enr}(\TPK)
\]
la troncature des histoires, et soit
\[
 q_{N,M}:\mathcal D_{\rad,N}^{\enr}\to
 \mathcal D_{\rad,M}^{\enr}
\]
l'élimination des données d'indice supérieur à \(M\).

\begin{theorem}[Naturalité du lecteur radial]
Le diagramme
\[
\begin{array}{ccc}
\Hist_N^{\enr}(\TPK)
 & \xrightarrow{\ \mathcal P_{\rad,N}^{\enr}\ }
 & \mathcal D_{\rad,N}^{\enr}\\[1.0em]
{\scriptstyle \pi_{N,M}}\big\downarrow
 && \big\downarrow{\scriptstyle q_{N,M}}\\[1.0em]
\Hist_M^{\enr}(\TPK)
 & \xrightarrow{\ \mathcal P_{\rad,M}^{\enr}\ }
 & \mathcal D_{\rad,M}^{\enr}
\end{array}
\]
commute ; c'est-à-dire,
\begin{equation}
 q_{N,M}\circ\mathcal P_{\rad,N}^{\enr}
 =\mathcal P_{\rad,M}^{\enr}\circ\pi_{N,M}.
 \label{espiral:eq:naturalidad-truncamiento}
\end{equation}
\end{theorem}
\begin{proof}
Les deux membres conservent exactement l'enregistrement initial et les \(M\) premières flèches, avec leurs cocycles locaux et accumulés, leurs facteurs affines, leurs holonomies préfixes et leurs entrées de registre de \(\widehat\gamma_N\). Aucune reconstruction à partir du produit final n'intervient : la séquence complète fait partie de la définition.
\end{proof}

\section{Lecteur de la limite inverse}

La compatibilité \cref{espiral:eq:naturalidad-truncamiento} induit
\begin{equation}
 \mathcal P_{\rad,\infty}^{\enr}:
 \varprojlim_N\Hist_N^{\enr}(\TPK)
 \longrightarrow
 \varprojlim_N\mathcal D_{\rad,N}^{\enr}.
 \label{espiral:eq:lector-limite-inverso}
\end{equation}
Chaque composante finie s'obtient par troncature de l'histoire complète. Le lecteur n'invente donc pas une continuité externe : il prolonge de manière compatible les certificats finis.

\section{Subdivision quadruple}

Dans le raffinement \(\sd_4\) du TPK, trois microflèches sont immobiles et la quatrième reproduit la transition observable. Si
\[
 \rho_{\rad}(e_1)=\rho_{\rad}(e_2)=\rho_{\rad}(e_3)=I,
 \qquad \rho_{\rad}(e_4)=h_e,
\]
alors
\[
 H_{\sd_4(e)}=I\,I\,I\,h_e=h_e.
\]
La naturalité est ainsi certifiée pour ce raffinement. Pour une substitution générale \(e\mapsto\widetilde e_k\cdots\widetilde e_1\), la condition exacte est
\[
 h_e=h_{\widetilde e_k}\cdots h_{\widetilde e_1}.
\]
Cette égalité, et non la ressemblance des courbes publiées, décide de la compatibilité d'un nouveau raffinement.

\section{Équivariance avec l'holonomie nonadique}

Soit \(b_9(x)\) le bloc de neuf transitions qui réalise \(\Gamma_9\) à partir de \(x\), et
\[
 H_9(x)=H_{b_9(x)}=(\Lambda_9(x),C_9(x)).
\]
L'action induite dans le registre radial ajoute ce bloc.

\begin{theorem}[Équivariance de cocycle]
On a
\begin{equation}
 \mathcal P_{\rad}^{\enr}\circ\Gamma_9
 =\widehat\Gamma_9^{\rad}\circ
 \mathcal P_{\rad}^{\enr}.
 \label{espiral:eq:equivarianza-gamma9}
\end{equation}
En composantes,
\[
 P\mapsto P+p(b_9),
 \qquad
 A\mapsto A+a(b_9),
\]
\[
 \Lambda\mapsto\Lambda_9(x)\Lambda,
 \qquad
 C\mapsto C_9(x)+\Lambda_9(x)C,
\]
\[
 g\mapsto g,\qquad q_{\mem}\mapsto q_{\mem}+1.
\]
\end{theorem}
\begin{proof}
L'action de \(\Gamma_9\) concatène \(b_9(x)\) à l'histoire, où \(p(b_9)\) et \(a(b_9)\) sont les sommes des caractères locaux de ses neuf flèches relevées. La loi de cocycle additionne les degrés et la loi affine \cref{espiral:eq:producto-affend} compose les holonomies. Le TPK apporte le retour de phase et l'incrément de mémoire. La définition de \(\widehat\Gamma_9^{\rad}\) reproduit exactement ces actualisations.
\end{proof}

Ce n'est pas une commutation littérale d'endomorphismes du même espace. Les codomaines sont distincts et la mémoire change. Le mot correct est \emph{équivariance}.

\section{Inversion dans le secteur réversible}

Si \(h_{e^\dagger}=h_e^{-1}\), le diagramme d'inversion satisfait
\[
 \mathcal P_{\rad}^{\enr}(C_{\mathrm{rev}}\gamma)
 =\mathcal J_{\rad}
 \mathcal P_{\rad}^{\enr}(\gamma),
\]
où \(\mathcal J_{\rad}\) inverse l'ordre des facteurs, remplace chaque facteur par son inverse et transporte le registre au moyen de l'involution TPK. Dans la sortie réelle, \cref{espiral:eq:lp-inversion} échange les deux orientations.

\chapter{Classification généalogique et non-fidélité de la publication}

\section{Signature radiale pleine d'une histoire}

Pour une histoire finie \(\widehat\gamma_N=\widehat e_N\cdots\widehat e_1\), nous définissons
\begin{equation}
 \mathfrak I_N(\widehat\gamma_N)=
 \left(
 (x_0,H_0,\Led_0);
 \bigl(\widehat e_j,p(\widehat e_j),a(\widehat e_j),P_j,A_j,
 h_{e_j},H_j,\Led_j\bigr)_{j=1}^{N}
 \right).
 \label{espiral:eq:invariante-espiral}
\end{equation}
La signature est pleine par rapport au lecteur radial enrichi choisi : elle conserve l'objet initial, la séquence typée de flèches, les caractères locaux et accumulés, tous les préfixes affines et le registre complet. Elle n'est pas présentée comme un invariant complet de toute l'HMT et ne remplace pas les coordonnées de l'état enrichi que le lecteur n'a pas pour domaine.

\begin{definition}[Équivalence spirale généalogique]
Deux histoires sont équivalentes s'il existe un isomorphisme typé de leurs données radiales qui transporte terme à terme la séquence de \cref{espiral:eq:invariante-espiral}, respecte les sources, les cibles, la composition et les registres, et commute avec toutes les troncatures déclarées.
\end{definition}

Dans la limite inverse, on obtient la famille compatible \((\mathfrak I_N)_{N\ge0}\). L'égalité des courbes dans \(\R^2\) est une relation plus grossière : l'évaluation archimédienne se factorise par l'équivalence généalogique, mais en général ne la reflète pas.

\section{Classification par 3 incréments visibles}

Soient \(\omega\) l'incrément angulaire, \(\beta\) l'incrément radial dans la coordonnée \(L_p\) et \(\mu\) l'incrément de mémoire relevée. Le tableau visible minimal est :
\begin{longtable}{ccccp{.34\textwidth}}
\caption{Classification cinématique des publications.}
\label{espiral:tab:clasificacion-cinematica}\\
\toprule
\(\omega\)&\(\beta\)&\(\mu\)&Type publié&Interprétation\\
\midrule
\endfirsthead
\toprule
\(\omega\)&\(\beta\)&\(\mu\)&Type publié&Interprétation\\
\midrule
\endhead
0&\(\neq0\)&0&recta radial&variation d'échelle sans rotation ni mémoire\\
\(\neq0\)&0&0&cercle&fermeture angulaire intrinsèquement radiale\\
\(\neq0\)&0&\(\neq0\)&hélice cylindrique&phase périodique avec mémoire relevée\\
\(\neq0\)&\(\neq0\)&0&espiral plana&variation radiale sans axe de mémoire visible\\
\(\neq0\)&\(\neq0\)&\(\neq0\)&suspension hélicoïdale radiale&rotation, auto-échelle et mémoire\\
0&0&\(\neq0\)&fibra axial&avancée de mémoire sans figure transverse\\
\bottomrule
\end{longtable}

Ce tableau classe la cinématique publiée, non toute la généalogie. Deux lignes ayant les trois mêmes incréments peuvent différer par \(p\), la retenue, le bord ou la classe affine.

\section{3 fibres d'une même évaluation visible}

\begin{example}[Cercle à mémoire nulle]
Soient \(H_j=I\), \(\mu=0\) et \(\theta_j=2\pi j/9\). La sortie est un cercle et le registre n'enregistre pas de relèvement axial.
\end{example}

\begin{example}[Le même cercle comme ombre]
Soient \(H_j=I\), \(\mu=1\) et la même phase angulaire. Dans \(\R^3\), on obtient l'hélice \cref{espiral:eq:inmersion-helice-nonadica} ; en appliquant \(\pi_{xy}\), la sortie plane coïncide avec celle de l'exemple précédent.
\end{example}

\begin{example}[Même spirale, retenue différente]
Soient deux histoires ayant les mêmes caractères réels \(p,\beta,\omega\), mais des décompositions de retenue distinctes que l'évaluation radiale envoie sur le même incrément \(\beta\). Toutes deux publient \(r_p(\theta)\), tandis que leurs registres diffèrent. Le lecteur enrichi les sépare ; la courbe, non.
\end{example}

\section{Théorème de classification généalogique}

\begin{theorem}[Correspondances spirales généalogiques]
Le module radial universel, la représentation locale \(e\mapsto h_e\) et les cocycles radiaux du relèvement TRIT étant fixés, toute histoire compatible TPK détermine :
\begin{enumerate}
\item une séquence naturelle de facteurs et d'holonomies affines préfixes ;
\item une classe généalogique stable sous troncature ;
\item une action équivariante de l'holonomie nonadique ;
\item dans le secteur réversible, une histoire conjuguée d'holonomie inverse ;
\item pour chaque secteur radial homogène admissible, une courbe régulière dans \(\R^2\) et, en conservant la profondeur, une suspension dans \(\R^3\).
\end{enumerate}
Une histoire de degré variable détermine en revanche un atlas ordonné de cartes par morceaux. L'application aux courbes n'est pas fidèle : elle peut identifier des histoires aux registres distincts.
\end{theorem}
\begin{proof}
Les points 1 et 2 sont les théorèmes d'holonomie affine et de naturalité. Le point 3 est \cref{espiral:eq:equivarianza-gamma9} ; le point 4 est l'inversion dans le sous-groupoïde réversible. Le point 5 s'obtient en restreignant la limite inverse au domaine \(\mathcal D_{\rad,\infty}^{\mathrm{hom}}\) et en appliquant ses caractères radial et angulaire. La non-fidélité découle de la \cref{espiral:fig:circulo-helice} et de la proposition générale du premier chapitre.
\end{proof}

\begin{statusbox}[title={Hypothèses, domaine et portée du théorème}]
L'holonomie nonadique, l'hélice, la vis de Pell et les cocycles du TPK sont des résultats récupérés. L'enveloppe affine, le lecteur par préfixes, sa naturalité et la publication \(L_p\) constituent une formalisation nouvelle exacte dans les domaines déclarés. La carte locale \(p(\widehat e)=\tau_+(e)+\tau_-(e)\) et \(a(\widehat e)=\tau_+(e)-\tau_-(e)\) est exacte dans la catégorie relevée de 2 demi-couronnes ; son utilisation hors de ce domaine exige de présenter le foncteur qui apporte les deux orientations. La publication sous forme d'une seule courbe est restreinte au secteur homogène admissible défini dans ce chapitre.
\end{statusbox}
